\section{Mina arbetsuppgifter under LIA-perioden}
\label{sec:arbetsuppgifter}

Under min LIA-period på Insighta, Inc arbetade jag med flera olika uppgifter inom DevOps och molninfrastruktur, vilket gav mig möjlighet att tillämpa mina teoretiska kunskaper i praktiken och samtidigt lära mig många nya verktyg.

Mitt första ticket handlade om att utvärdera om K8sRunLauncher kunde användas för att deploya Dagster OSS utan att avbryta pågående jobb. Genom detta arbete lärde jag mig mycket om agila arbetsmetoder, som vi fått en introduktion till på yrkeshögskolan, samt om pull requests, GitHub Actions och arbetsflöden. Jag fick även praktisk erfarenhet av Kubernetes, Docker och Docker-compose, och jag lärde mig hur man planerar, testar och dokumenterar lösningar tillsammans med teamet.

Ett annat viktigt ticket var att flytta Dagster OSS och Postgres från Docker till Kubernetes. Tidigare kördes Dagster på virtuella maskiner, men vi ville samla all komputering i samma system för att göra det enklare och mer skalbart. Jag lärde mig hur man deployar Dagster på Kubernetes, konfigurerar så att jobben skapas automatiskt, kopplar till Postgres och sparar loggar i en GCP-bucket. Detta gav mig en mycket bra förståelse för hur man bygger fungerande pipelines och molninfrastruktur.

Slutligen arbetade jag med att göra branch-deploys med cloud-baserade Kubernetes-kluster, framför allt på GKE. Tidigare använde vi separata virtuella maskiner för varje branch, men nu ville vi flytta deploymenten till cloud-kluster för att kunna skala, isolera och automatisera bättre. Här lärde jag mig hur man sätter upp GKE-kluster, hanterar resurser och isolering för varje branch, och använder Helm charts för att hantera deploymenten. Dessutom fick jag första erfarenheten av Terraform CDKTF, något som inte ingick i våra kurser på yrkeshögskolan. Jag dokumenterade stegen och lärdomarna, vilket hjälpte teamet att förstå och förbättra branch-deploy-strategin i molnet.

Sammantaget har dessa uppgifter gett mig erfarenhet av att arbeta med verkliga DevOps-projekt, använda moderna verktyg och molninfrastruktur, samt att samarbeta med ett team enligt agila principer. Jag har fått praktisk erfarenhet av infrastruktur som kod, CI/CD, Kubernetes, GCP och Dagster, och jag har också utvecklat min förmåga att dokumentera och dela lärdomar med kollegor.